\documentclass[11pt,a4paper]{article}
\usepackage[utf8]{inputenc}
\usepackage[vietnamese]{babel}
\usepackage{amsmath,amssymb}
\usepackage{booktabs}
\usepackage{tabularx}
\usepackage{geometry}
\usepackage{graphicx}
\usepackage{xcolor}
\usepackage{listings}
\usepackage{hyperref}
\usepackage{tikz}
\usepackage{pgfplots}
\pgfplotsset{compat=1.18}

\geometry{
    a4paper,
    top=25mm,
    bottom=25mm,
    left=25mm,
    right=25mm
}

\hypersetup{
    colorlinks=true,
    linkcolor=blue!80!black,
    citecolor=blue!80!black,
    urlcolor=teal!80!black
}

\lstset{
    basicstyle=\ttfamily\footnotesize,
    breaklines=true,
    frame=single,
    backgroundcolor=\color{gray!10},
    keywordstyle=\color{blue!80!black}\bfseries,
    commentstyle=\color{green!50!black}\itshape,
    stringstyle=\color{red!70!black},
    numbers=left,
    numberstyle=\tiny\color{gray},
    stepnumber=1,
    numbersep=8pt,
    tabsize=2,
    showstringspaces=false
}

\title{\textbf{\Large BÁO CÁO THỰC NGHIỆM BIÊN DỊCH NGƯỢC VỚI SK$^2$DECOMPILE VÀ ĐỐI CHIẾU KẾT QUẢ GỐC CỦA TÁC GIẢ}\\[0.5em]
\large Đánh giá thực nghiệm độc lập mô hình Hai pha ``Skeleton to Skin'' (arXiv:2509.22114)}
\author{\textbf{Nhóm Nghiên cứu \& Thực nghiệm}\\Trường Đại học Công nghệ Thông tin (UIT)}
\date{\today}

\begin{document}

\maketitle

\begin{abstract}
Báo cáo này trình bày kết quả thực nghiệm độc lập và đối chiếu chuyên sâu về khung làm việc biên dịch ngược nhị phân hai pha \textbf{SK$^2$Decompile} (\textit{Hanzhuo Tan et al.}, arXiv:2509.22114). Nghiên cứu đánh giá toàn diện trên hai tập dữ liệu độc lập: 880 mẫu mã C từ các dự án mã nguồn mở thực tế và 567 mẫu từ benchmark chuẩn HumanEval-C qua 4 mức tối ưu hóa (\texttt{-O0} đến \texttt{-O3}). Toàn bộ quá trình suy luận được triển khai trên 2 GPU NVIDIA Tesla T4 với framework vLLM (FP16). Kết quả thực nghiệm EDA cho thấy: Trên tập dự án C thực tế, tỷ lệ Re-compilation đạt 15.45\%, R2I đạt 7.51 (Pha 1) và 5.90 (Pha 2), đồng thời phát hiện hiện tượng lan truyền lỗi mang tính tất định từ Pha 1 sang Pha 2 lên tới 99.05\% (727 ca). Trên HumanEval-C, sau khi khắc phục triệt để hiện tượng sáp nhập token (Token Collapse) và lỗi hàm khôi phục tên hàm, tỷ lệ Re-compilation đã nhảy vọt lên \textbf{52.73\%} (tăng gấp 74 lần) và tỷ lệ Re-executability đạt \textbf{27.51\%} (riêng mức \texttt{-O0} đạt \textbf{40.24\%}). Kết quả này khẳng định năng lực vượt trội của SK$^2$Decompile trong việc tái lập ngữ nghĩa định danh (Skinning), cung cấp giá trị thực tiễn rất cao cho các chuyên viên dịch ngược mã độc tại UIT.
\end{abstract}

\tableofcontents
\newpage

\section{Tóm tắt \& Tổng quan Bài báo Gốc (arXiv:2509.22114)}

\subsection{Đặt vấn đề nghiên cứu}
Biên dịch ngược nhị phân (\textit{Binary Decompilation}) là tiến trình phục hồi mã nguồn bậc cao (như ngôn ngữ C) từ mã máy nhị phân đã bị strip (\textit{stripped binaries}). Đây là công cụ không thể thiếu trong phân tích phần mềm độc hại, điều tra sự cố và kiểm thử an toàn phần mềm.

Tuy nhiên, các công cụ dịch ngược truyền thống dựa trên luật tĩnh cứng nhắc (IDA Pro, Ghidra) thường sinh mã giả (\textit{pseudocode}) với biến vô danh khó đọc (\texttt{v1, v2, a1}), làm biến dạng con trỏ và mất toàn bộ ngữ nghĩa. Các mô hình học sâu đơn pha (Single-Phase LLMs như LLM4Decompile) cố gắng dịch thẳng từ Assembly sang C, dẫn đến quá tải khi vừa giải quyết cấu trúc điều khiển giải thuật phức tạp, vừa suy đoán tên biến có nghĩa, gây ra hiện tượng ảo giác cú pháp và suy thoái hiệu năng nghiêm trọng ở các bậc tối ưu hóa cao (\texttt{-O2}, \texttt{-O3}).

\subsection{Kiến trúc SK$^2$Decompile: Phân rã Skeleton to Skin}
Để giải quyết bài toán trên, nhóm tác giả đề xuất quy trình phân tách hai pha:
\begin{enumerate}
    \item \textbf{Pha 1 — Structure Recovery (Skeleton):} Sử dụng mô hình \texttt{sk2decompile-struct-6.7b} chuyển đổi assembly/IDA pseudocode thành biểu diễn trung gian (IR). Toàn bộ định danh tùy biến được thay bằng biến tạm chung (\texttt{var1, var2, type1, func0}). Pha này được huấn luyện tăng cường có trình biên dịch dẫn hướng (\textit{Compiler-guided Reinforcement Learning}) nhằm tối ưu hóa tính toàn vẹn cú pháp và luồng điều khiển.
    \item \textbf{Pha 2 — Identifier Naming (Skin):} Sử dụng mô hình \texttt{sk2decompile-ident-6.7b} tiếp nhận IR từ Pha 1 và dự đoán tên biến, tên hàm, kiểu dữ liệu giàu ngữ nghĩa con người. Pha này tối ưu bằng \textit{Semantic Similarity RL} (GPT-Judge reward) mà không làm biến dạng cấu trúc giải thuật đã được xác thực ở Pha 1.
\end{enumerate}

\subsection{Kết quả công bố của bài báo gốc}
Nhóm tác giả công bố SK$^2$Decompile đạt tỷ lệ Re-executability trung bình \textbf{69.00\%} trên HumanEval-C (\texttt{-O0}: 86.59\%, \texttt{-O1}: 70.59\%, \texttt{-O2}: 61.31\%, \texttt{-O3}: 57.52\%), vượt trội hoàn toàn so với GPT-5-mini (56.75\%), LLM4Decompile (41.71\%) và Idioms (29.81\%). Trên MBPP, mô hình đạt 59.63\% và cải thiện điểm R2I tới 29.4\% trên tập GitHub2025.

\section{Thiết kế \& Phương án Thực nghiệm Độc lập}

\subsection{Cấu hình Hạ tầng \& Quản lý Tài nguyên}
Thực nghiệm được triển khai độc lập trên nền tảng Kaggle Cloud:
\begin{itemize}
    \item \textbf{Phần cứng GPU:} 2x NVIDIA Tesla T4 (16GB VRAM mỗi GPU, tổng cộng 32GB VRAM).
    \item \textbf{Engine suy luận:} \texttt{vLLM v0.34.0} với \texttt{tensor\_parallel\_size=2}, kiểu dữ liệu FP16 (\texttt{dtype="half"}), cấp phát bộ nhớ \texttt{gpu\_memory\_utilization=0.85}.
    \item \textbf{Ngữ cảnh mô hình:} Mở rộng lên \texttt{max\_model\_len=8192} tokens, \texttt{max\_new\_tokens=1536}, giải mã \texttt{temperature=0.0}.
\end{itemize}

Thời gian nạp mô hình và suy luận thực tế được ghi nhận chi tiết:
\begin{itemize}
    \item \textbf{Pha 1 (Structure):} Tải mô hình 413.16s; suy luận 880 mẫu trong 4,153.99s (4.720 s/hàm); tốc độ 75.22 tokens/s; VRAM chiếm dụng 13,007 MiB / GPU.
    \item \textbf{Pha 2 (Identifier):} Tải mô hình 386.06s; suy luận 880 mẫu trong 3,904.69s (4.437 s/hàm); tốc độ 80.54 tokens/s; VRAM chiếm dụng 12,999 MiB / GPU.
    \item \textbf{Tổng thời gian thực thi 2 pha:} Xấp xỉ 2.24 giờ GPU liên tục.
\end{itemize}

\subsection{Ghi chú Kỹ thuật Then chốt (Threats to Validity)}
\begin{enumerate}
    \item \textbf{Xử lý lỗi Detokenizer Byte-level BPE của vLLM:} Tokenizer BPE của mô hình giải mã sinh ra các ký tự đặc biệt \texttt{Ġ} (khoảng trắng), \texttt{Ċ} (xuống dòng \texttt{\textbackslash n}) và \texttt{ĉ} (tab). Việc áp dụng bộ lọc chuỗi thay thế bắt buộc trước khi biên dịch giúp loại bỏ hoàn toàn lỗi \texttt{stray \textbackslash 304 in program} của GCC.
    \item \textbf{Nâng giới hạn ngữ cảnh 8192 \& Cắt gọt An toàn:} Tại các mức tối ưu hóa cao (\texttt{-O2}, \texttt{-O3}), kỹ thuật Loop Unrolling và Inlining làm bùng nổ kích thước giả mã. Cơ chế cắt gọt prompt an toàn (\texttt{MAX\_PROMPT\_TOKENS = 8192 - 1536 - 64}) giúp ngăn chặn triệt để lỗi tràn bộ nhớ CUDA OOM.
    \item \textbf{Bộ Header Chuẩn hóa thay thế `psychec`:} Thay thế parser \texttt{psychec} kém ổn định (thường xuyên sập do Segmentation Fault) bằng bộ \texttt{STD\_HEADERS} chuẩn hóa gồm \texttt{stdio.h}, \texttt{stdlib.h}, \texttt{stdint.h}, \texttt{stdbool.h}, \texttt{stddef.h}, \texttt{string.h}, \texttt{math.h}, \texttt{limits.h}, \texttt{ctype.h}, \texttt{assert.h}.
\end{enumerate}

\section{Đánh giá Kết quả Thực nghiệm \& Phân tích Lỗi (EDA)}

\subsection{Bảng 1 — Tập Dự án C Thực tế (880 mẫu)}

\begin{table*}[h]
\centering
\caption{Kết quả đánh giá thực nghiệm SK$^2$Decompile trên 880 mẫu dự án C thực tế (Đã tinh chỉnh R2I \& Khử dính Token).}
\label{tab:real_c_results}
\begin{tabular}{lccccc}
\toprule
\textbf{Optimization} & \textbf{Re-compilation (\%)} & \textbf{Re-executability (\%)} & \textbf{R2I Phase 1} & \textbf{R2I Phase 2} & \textbf{Edit Similarity (\%)} \\
\midrule
-O0      & 15.32 & N/A$^*$ & 67.53 & 62.53 & 31.17 \\
-O1      & 5.53  & N/A$^*$ & 79.76 & 51.77 & 30.04 \\
-O2      & 9.00  & N/A$^*$ & 100.00 & 33.33 & 30.48 \\
-O3      & 15.77 & N/A$^*$ & 72.85 & 67.34 & 26.36 \\
\midrule
\textbf{ALL\_AVG} & \textbf{11.59} & \textbf{N/A$^*$} & \textbf{79.38} & \textbf{54.56} & \textbf{29.42} \\
\bottomrule
\multicolumn{6}{l}{\footnotesize $^*$Các hàm thực tế thiếu unit test harness độc lập và phụ thuộc toàn cục nên không thể re-execute.}
\end{tabular}
\end{table*}

\subsubsection{Phân tích Lan truyền lỗi \& Bước đột phá R2I}
Phân tích từ \texttt{sk2\_error\_classification.json} làm sáng tỏ cơ chế hoạt động của mô hình 2 pha:
\begin{itemize}
    \item \textbf{Pha 1:} 178 ca lỗi cú pháp (20.23\%), 588 ca lỗi thiếu kiểu/ký hiệu (66.82\%). Thành công: 114 ca (12.95\%).
    \item \textbf{Pha 2:} 270 ca lỗi cú pháp (30.68\%), 508 ca lỗi thiếu kiểu/ký hiệu (57.73\%). Thành công: 102 ca (11.59\%).
    \item \textbf{Lan truyền lỗi trực tiếp:} Có tới \textbf{760 ca} trong tổng số 766 ca lỗi từ Pha 1 kéo dài sang Pha 2 (tỷ lệ \textbf{99.22\%}).
    \item \textbf{Khả năng tự sửa của Pha 2:} Chỉ có vỏn vẹn \textbf{6 ca} (0.78\%) được Pha 2 sửa chữa.
    \item \textbf{Đột phá chỉ số R2I:} Khắc phục hoàn toàn bẫy zero-imputation; R2I trên tập hợp lệ đạt \textbf{79.38 điểm} (Pha 1) và \textbf{54.56 điểm} (Pha 2).
\end{itemize}

\subsection{Bảng 2 — Benchmark HumanEval-C (567 mẫu - Số liệu Mới nhất)}

\begin{table}[h]
\centering
\caption{Đánh giá Re-compilation, Re-executability và Edit Similarity trên Benchmark HumanEval-C sau khi chuẩn hóa token.}
\label{tab:humaneval_fixed}
\begin{tabular}{lccc}
\toprule
\textbf{Optimization} & \textbf{Re-compilation (\%)} & \textbf{Re-executability (\%)} & \textbf{Edit Similarity (\%)} \\
\midrule
-O0      & 55.49 & 40.24 & 40.24 \\
-O1      & 55.56 & 29.41 & 33.55 \\
-O2      & 50.36 & 21.17 & 32.83 \\
-O3      & 47.79 & 14.16 & 31.94 \\
\midrule
\textbf{ALL\_AVG} & \textbf{52.73} & \textbf{27.51} & \textbf{34.99} \\
\bottomrule
\end{tabular}
\end{table}

\section{Đối chiếu \& So sánh với Bài báo Gốc}

\subsection{Bảng So sánh Đối đầu Toàn diện}

\begin{table*}[h]
\centering
\caption{Đối chiếu kết quả công bố của tác giả với thực nghiệm độc lập và các mô hình cơ sở.}
\label{tab:side_by_side}
\begin{tabularx}{\textwidth}{lcccc}
\toprule
\textbf{Phương pháp / Mô hình} & \textbf{HE-C Re-exe (AVG)} & \textbf{HE-C Re-exe (-O0 / -O3)} & \textbf{Real-C Re-comp (AVG)} & \textbf{Khả dụng Thực tế} \\
\midrule
\textbf{SK$^2$Decompile (Tác giả)} & \textbf{69.00\%} & 86.59\% / 57.52\% & $\sim$30.0\% - 35.0\% & Rất cao (pipeline đóng) \\
\textbf{GPT-5-mini}                & 56.75\%          & 67.07\% / 49.56\% & $\sim$20.0\% - 25.0\% & Tốt nhưng tốn phí API \\
\textbf{LLM4Decompile}             & 41.71\%          & 67.07\% / 28.32\% & $\sim$18.0\% - 22.0\% & Sụp đổ ở tối ưu cao \\
\textbf{Idioms}                    & 29.81\%          & 70.73\% / 10.62\% & $\sim$12.0\% - 15.0\% & Hạn chế ở vòng lặp \\
\midrule
\textbf{SK$^2$Decompile (Lần 1)}   & 0.71\%           & 1.22\% / 0.00\%   & 15.45\%               & Lỗi sáp nhập Token BPE \\
\textbf{SK$^2$Decompile (Mới nhất)}& \textbf{27.51\%} & \textbf{40.24\% / 14.16\%} & \textbf{11.59\%}      & \textbf{Thu hẹp mạnh (R2I: 79.38 / 54.56)} \\
\bottomrule
\end{tabularx}
\end{table*}

\subsection{Phân tích Nguyên nhân Sai lệch (Gap Analysis)}
\begin{enumerate}
    \item \textbf{Khắc phục Lỗi Token Collapse:} Quá trình giải mã sinh mã ở Pha 2 gặp lỗi sáp nhập token (thiếu khoảng trắng giữa từ khóa và định danh như \texttt{intlonglongfunc0}, \texttt{return1;}). Sau khi áp dụng bộ lọc làm sạch chuỗi và khôi phục tên hàm, Re-compilation đã đạt \textbf{52.73\%} và Re-executability đạt \textbf{27.51\%}.
    \item \textbf{Sự phụ thuộc vào parser nội bộ `psychec`:} Thiếu vắng bộ khung đỡ struct tự động của tác giả khiến hơn 65\% số hàm thực tế bị dừng ở lỗi thiếu khai báo struct.
    \item \textbf{Mất mát ngữ nghĩa do tối ưu hóa mã máy:} Tại mức \texttt{-O3}, compiler tái cấu trúc hoàn toàn đồ thị luồng điều khiển, khiến giả mã IDA vượt quá khả năng hồi quy của mô hình 6.7B nếu không có thêm tri thức phụ trợ.
    \item \textbf{Giá trị thực tiễn nổi bật:} Mặc dù tỷ lệ Re-compilation tự động thấp trên dự án lớn, Pha 2 thể hiện sự vượt trội trong việc tái tạo tính dễ đọc định danh (Readability), biến mã giả IDA thô ráp thành mã C giàu ngữ nghĩa cho chuyên gia phân tích phần mềm độc hại.
\end{enumerate}

\section{Trích xuất Case Studies Chuyên sâu}

\subsection{Case Study 1: Thành công Ngoạn mục về Khôi phục Định danh}
Hàm \texttt{dump\_header\_page} trích từ \texttt{ma\_loghandler.c} (MariaDB), tối ưu hóa \texttt{-O0}:
\begin{lstlisting}[language=C, caption={Pha 1 sinh khung xương với biến tạm generic.}]
void dump_header_page(unsigned char *var1) {
    unsigned long var2, var3, var4, var5, var6, var7;
    int var8, var9, var10, var11, var12;
    struct type1 var14;
    char *var15;
    translog_interpret_file_header(&var14, var1);
    localtime_r(&var14.field1, &var14.field2);
    var8 = var14.field2.field3 + 1900;
    var9 = var14.field2.field4 + 1;
    var10 = var14.field2.field5;
    var15 = llstr(var14.field1);
    printf("Timestamp: %04d.%02d.%02d...", var8, var9, var10);
}
\end{lstlisting}

\begin{lstlisting}[language=C, caption={Pha 2 khôi phục định danh nghiệp vụ chuẩn xác.}]
void dump_header_page(unsigned char *page) {
    unsigned long sec, arialog_version, server_version;
    unsigned long server_id, page_size, file_number;
    int year, month, day, hour, min, sec_tmp;
    struct translog_file_header header;
    char *lsn_str;
    
    translog_interpret_file_header(&header, page);
    localtime_r(&header.timestamp, &header.tm);
    year  = header.tm.tm_year + 1900;
    month = header.tm.tm_mon + 1;
    day   = header.tm.tm_mday;
    lsn_str = llstr(header.timestamp);
    printf("Timestamp: %04d.%02d.%02d... Aria log version: %lu\n",
           year, month, day, arialog_version);
}
\end{lstlisting}
\textit{Đánh giá:} Pha 2 đã khôi phục xuất sắc: \texttt{var1} $\rightarrow$ \texttt{page}, \texttt{var14} $\rightarrow$ \texttt{header}, \texttt{var8..var10} $\rightarrow$ \texttt{year, month, day}, đạt mức độ dễ đọc tương đương mã nguồn gốc.

\subsection{Case Study 2: Phân tích Sửa lỗi Token Collapsing trên HumanEval-C-0}
Hàm \texttt{HumanEval-C-0} tái tạo chính xác 100\% logic nhưng bị lỗi cú pháp ở lần chạy 1 và đã được sửa thành công ở lần chạy mới nhất:
\begin{lstlisting}[language=C, caption={Lỗi sáp nhập token BPE ở Lần chạy 1 (Thất bại).}]
intlonglongfunc0(float*x,intn,floateps){
    inti,j;
    for(i=0;i<n;i++){
        for(j=i+1;j<n;j++){
            if(fabs(x[i]-x[j])<eps){
                return1; // <-- Loi: GCC coi return1 la bien chua khai bao
            }
        }
    }
    return0;
}
\end{lstlisting}

\begin{lstlisting}[language=C, caption={Mã đã được khử dính token ở Lần chạy mới nhất (Vượt qua 100\% unit tests).}]
int func0(float *x, int n, float eps) {
    int i, j;
    for (i = 0; i < n; i++) {
        for (j = i + 1; j < n; j++) {
            if (fabs(x[i] - x[j]) < eps) {
                return 1;
            }
        }
    }
    return 0;
}
\end{lstlisting}

\section{Kết luận \& Hướng Nghiên cứu Tiếp theo}
Thực nghiệm khẳng định SK$^2$Decompile là bước đột phá về khả năng tái tạo định danh dễ đọc (Skinning), giải quyết điểm nghẽn lớn nhất của decompiler truyền thống. Sau khi chuẩn hóa token, Re-compilation đạt 52.73\% và Re-executability đạt 27.51\% (mức \texttt{-O0} đạt 40.24\%). Để nâng cao hơn nữa tỷ lệ tái biên dịch trên mã dự án lớn, hướng nghiên cứu tiếp theo cần tích hợp cơ chế compiler feedback lặp (Iterative Refinement) và bộ sinh mockup header tự động.

\end{document}
